\documentclass[a4paper,11pt]{article}

\usepackage[utf8]{inputenc}
\usepackage{amsmath, amssymb}
\usepackage{graphicx}
\usepackage{algorithm2e}

\title{L'algoritme \textit{k-means}}
\author{}
\date{}

\begin{document}

\maketitle

\section{Introducció}

L'algoritme \textit{k-means} és un dels mètodes de \emph{clustering}
(no supervisat) més utilitzats.
Donat un conjunt de dades, el seu objectiu és agrupar els punts en $k$
grups o clústers de manera que els punts d'un mateix clúster siguin
més semblants entre ells que als d'altres clústers.

Intuïtivament, \textit{k-means} busca $k$ centres (anomenats \emph{centroides})
de manera que cada punt de dades s’assigni al centre més proper,
i a continuació actualitza aquests centres com la mitjana dels punts
assignats. Aquest procés es repeteix iterativament fins que el sistema
``s'estabilitza''.

\section{Formulació matemàtica}

Sigui un conjunt de dades
\[
  X = \{x_1, x_2, \dots, x_n\}, \quad x_i \in \mathbb{R}^d,
\]
i un nombre de clústers $k \in \mathbb{N}$ tal que $1 \leq k < n$.

L'objectiu de \textit{k-means} és trobar:

\begin{itemize}
  \item Una partició de les dades en $k$ clústers $C_1, C_2, \dots, C_k$
        tals que $C_j \subset \{1,\dots,n\}$, $C_j \neq \emptyset$,
        $C_j \cap C_\ell = \emptyset$ si $j \neq \ell$, i 
        $\bigcup_{j=1}^k C_j = \{1,\dots,n\}$. Informalment, $C_i$ es
		un subconjunt no buit de $\{1,\dots,n\}$ on les dades
		dels indexos pertanyents a $C_i$ son semblants.
  \item Un conjunt de centroides
        \[
          \mu_1, \mu_2, \dots, \mu_k \in \mathbb{R}^d.
        \]
	Que representa el centre de cada cluster.
\end{itemize}

L'objectiu del problema es separar les dades en K clusters, de manera que
una dada pertany a un cluster i que les dades dins de cada cluster siguin properes
entre elles. El problema es planteja com la minimització de la suma d'errors
quadràtics dins de cada clúster (funció objectiu de \textit{k-means}):

\[
  J(\mu_1,\dots,\mu_k, C_1,\dots,C_k)
  = \sum_{j=1}^{k} \sum_{i \in C_j} \| x_i - \mu_j \|^2.
\]

La resolució d'aquest problema de minimització no és trivial. Kmeans és un
algoritme iteratiu que intenta convergir al mínim. L'algoritme \textit{k-means} 
clàssic (també conegut com a algoritme de Lloyd) alterna dos passos:

\begin{enumerate}
  \item \textbf{Pas d'assignació}.  
  Donat un conjunt de centroides $\mu_1, \mu_2, \dots, \mu_k \in \mathbb{R}^d$, el conjunt de punts assignats al clúster $j$ és
  \[
    C_j = \left\{\, i \;\middle|\; 
      \|x_i - \mu_j\|^2 \leq \|x_i - \mu_\ell\|^2 \;\; \forall \ell \in \{1,\dots,k\}
    \right\}.
  \]
  És a dir, cada punt s’assigna al centre més proper.

  \item \textbf{Pas d’actualització}.  
  Un cop fetes les assignacions, cada centre s’actualitza com la mitjana dels punts del seu clúster:
  \[
    \mu_j = \frac{1}{|C_j|} \sum_{i \in C_j} x_i,
    \qquad j = 1,\dots,k.
  \]
\end{enumerate}

Aquests dos passos es repeteixen fins que l’assignació deixa de canviar
o fins assolir un nombre màxim d'iteracions.

\section{Descripció detallada de l'algoritme}

\subsection{Inicialització}

La primera decisió que s'ha de prendre és com escollir els centroides
inicials $\mu_1^{(0)}, \dots, \mu_k^{(0)}$.

Opcions típiques:
\begin{itemize}
  \item Escollir aleatòriament $k$ punts del conjunt de dades i usar-los
        com a centroides inicials.
  \item Utilitzar heurístiques com \textit{k-means++}, que escull
        els centroides inicials amb una probabilitat proporcional
        a la distància al centre més proper ja seleccionat.
\end{itemize}

La inicialització és important perquè \textit{k-means} pot convergir
a mínims locals: diferents inicialitzacions poden donar lloc a solucions
diferents.

\subsection{Pas d’assignació (E-step)}

Donat un conjunt de centroides fixos
$\mu_1, \dots, \mu_k$, l’algoritme recorre cada punt $x_i$
i calcula les seves distàncies a cada centre. S’assigna $x_i$ al clúster
que tingui el centre més proper.

Des del punt de vista de la funció objectiu $J$, aquest pas minimitza
$J$ respecte a les variables d’assignació (els $C_j$) mantenint fixos
els centres.

\subsection{Pas d’actualització (M-step)}

Un cop assignats els punts als clústers, l'algoritme recalcula
els centroides com la mitjana dels punts assignats a cada clúster:

\[
  \mu_j = \frac{1}{|C_j|} \sum_{i \in C_j} x_i.
\]

Aquest pas minimitza la funció objectiu $J$ respecte als centroides,
mantenint fixa la partició dels punts.

\subsection{Criteris d’aturada}

L'algoritme es repeteix fins que es compleix alguna condició:

\begin{itemize}
  \item Les assignacions dels punts a clústers deixen de canviar.
  \item La variació dels centroides entre iteracions consecutives
        és menor que un llindar
        \[
          \max_j \|\mu_j^{(t)} - \mu_j^{(t-1)}\| < \varepsilon.
        \]
  \item S'assoleix un nombre màxim d'iteracions $T_{\max}$.
\end{itemize}

\section{Pseudocodi de l'algoritme \textit{k-means}}

A continuació presentem un pseudocodi estàndard per a l'algoritme
\textit{k-means}. Assumim que la funció
\texttt{distancia(x, y)} calcula la distància (per exemple, euclidiana)
entre els vectors \texttt{x} i \texttt{y}.

\section{Pseudocodi de l'algoritme \textit{k-means}}

\begin{algorithm}[H]
\DontPrintSemicolon
\caption{Algoritme \textit{k-means}}

\KwIn{Conjunt de dades $X = \{x_1,\dots,x_n\}$, nombre de clústers $k$, llindar $\varepsilon$, iteracions màximes $T_{\max}$.}
\KwOut{Clústers $C_1,\dots,C_k$ i centroides $\mu_1,\dots,\mu_k$.}

\BlankLine
\textbf{Inicialització:}\;
Seleccionar aleatòriament $k$ punts de $X$ i definir-los com a centroides inicials $\mu_1,\dots,\mu_k$.\;
$t \leftarrow 0$.\;

\Repeat(\tcp*[f]{Criteri d'aturada}){ $\Delta < \varepsilon$ \textbf{o} $t \ge T_{\max}$ }{

  \BlankLine
  \tcp{Pas d'assignació}
  \For{$i = 1$ \KwTo $n$}{
    Assignar $x_i$ al clúster:
    \[
      j^\star = \arg\min_{1 \le j \le k} \|x_i - \mu_j\|^2
    \]
    \[
      x_i \in C_{j^\star}
    \]\;
  }

  \BlankLine
  \tcp{Guardar antics centroides}
  $\mu_j^{\text{antic}} \leftarrow \mu_j$ per a tot $j$.\;

  \BlankLine
  \tcp{Pas d'actualització}
  \For{$j = 1$ \KwTo $k$}{
    Recalcular el centroide:
    \[
      \mu_j \leftarrow \frac{1}{|C_j|} \sum_{i \in C_j} x_i
    \]\;
  }

  \BlankLine
  \tcp{Càlcul del desplaçament per al criteri de parada}
  $\Delta \leftarrow \max_{1 \le j \le k} \|\mu_j - \mu_j^{\text{antic}}\|$\;

  $t \leftarrow t + 1$\;

}

\Return{els clústers $C_1,\dots,C_k$ i centroides finals $\mu_1,\dots,\mu_k$}\;

\end{algorithm}

\section{Comentaris i consideracions addicionals}

\subsection{Complexitat}

A cada iteració, el pas d'assignació requereix calcular distàncies
de cada punt a cada centroide, amb una complexitat
$O(nkd)$ si la dimensió és $d$.
El pas d’actualització també és $O(nd)$ (només cal recorrer un cop els punts
per calcular les mitjanes). Per tant,
si l’algoritme executa $T$ iteracions, la complexitat total és
aproximadament
\[
  O(T \cdot n \cdot k \cdot d).
\]

\subsection{Sensibilitat a la inicialització}

\textit{k-means} és sensible a com s’inicialitzen els centroides:
diferents inicialitzacions poden donar solucions de qualitat diferent.
Per això sovint s’executa diverses vegades amb diferents llavors
i s’escull la solució amb menor valor de $J$.

\subsection{Limitacions}

Algunes limitacions clàssiques de \textit{k-means}:

\begin{itemize}
  \item Assumeix que els clústers són aproximadament esfèrics
        (en l’espai de la mètrica utilitzada).
  \item No gestiona bé clústers amb mides molt diferents
        o densitats desiguals.
  \item És sensible a la presència d’\emph{outliers}, ja que
        utilitza la mitjana per calcular els centroides.
  \item Requereix fixar $k$ prèviament; escollir el nombre adequat
        de clústers no sempre és trivial.
\end{itemize}

\section{Adaptació de l'algoritme \textit{k-means} al cas de les enquestes}

En el context del projecte, cada persona respon un conjunt de preguntes
diferents, cadascuna de les quals pot ser de tipologia distinta.
Per tal d’aplicar mètodes de \textit{clustering}, representarem cadascuna
d’aquestes respostes com un vector en un espai multidimensional.
Així doncs, cada individu queda representat per un vector
\[
  X_i = [X_{i1}, X_{i2}, \dots, X_{iN}],
\]
on $X_{ik}$ és la resposta de la persona $i$ a la pregunta $k$.

El conjunt de totes les respostes de totes les persones forma una matriu
\[
  X =
  \begin{bmatrix}
    X_{11} & \cdots & X_{1N} \\
    \vdots & \ddots & \vdots \\
    X_{M1} & \cdots & X_{MN}
  \end{bmatrix},
\]
on $M$ és el nombre de persones i $N$ el nombre total de preguntes.

\subsection{Tipologies de variables}

Cada component $X_{ik}$ del vector pot ser de tipologia diferent,
fet que introdueix la necessitat d’una funció de distància heterogènia.
Les tipologies considerades són:

\begin{itemize}
  \item \textbf{Variables quantitatives (numèriques):}
        emmagatzemen un valor numèric real.
  \item \textbf{Variables qualitatives ordenades:}
        un sol valor pertanyent a un conjunt de modalitats ordenades
        (p. ex.\ \textit{poc, normal, força, molt}).
  \item \textbf{Variables qualitatives no ordenades:}
        un únic valor d’un conjunt de modalitats sense ordre.
  \item \textbf{Variables qualitatives no ordenades multivaluades:}
        una resposta formada per un conjunt de $p$ modalitats
        possibles d'entre un conjunt de $q$ opcions.
  \item \textbf{Respostes de text lliure:}
        un string arbitrari sense valors predefinits.
\end{itemize}

Aquestes diferències impedeixen usar directament distàncies clàssiques
com la euclidiana sense una definició local de distància adaptada.

\subsection{Distàncies locals i global}

La distància entre dues persones $X_i$ i $X_j$ s’obté com la mitjana
de les distàncies locals:
\[
  D(X_i , X_j) = \frac{1}{N}\sum_{k=1}^N D_\text{local}(X_{ik}, X_{jk}),
\]
on $D_\text{local}$ depèn de la tipologia de la variable:

\begin{enumerate}
  \item \textbf{Numèrica:}
        \[
          D_\text{local}(x,y)=\frac{|x-y|}{V_{\max}-V_{\min}}.
        \]
  \item \textbf{Qualitativa ordenada:}
        \[
          D_\text{local}(x,y)=\frac{|m(x)-m(y)|}{\#\text{mod}-1}.
        \]
  \item \textbf{Qualitativa no ordenada:}
        \[
          D_\text{local}(x,y)=
          \begin{cases}
            0 & x=y,\\
            1 & x\neq y.
          \end{cases}
        \]
  \item \textbf{Multivaluada (Jaccard):}
        \[
          D_\text{local}(A,B)=1-\frac{|A\cap B|}{|A\cup B|}.
        \]
  \item \textbf{Text lliure (Levenshtein):}
	  \[
		  D_\text{local}(s,t) = \frac{L_{s,t}(\text{len}(s), \text{len}(t)) - \left| \text{len}(s) - \text{len}(t) \right|}{\max(\text{len}(s), \text{len}(t)) - \left| \text{len}(s) - \text{len}(t) \right|}
	\]
        on $L(s,t)$ és la \textbf{distància de Levenshtein}.
\end{enumerate}

\subsubsection*{Definició de la distància de Levenshtein}

Donats dos strings \( s \) i \( t \), la distància de Levenshtein \( L(s,t) \) és:

\[
L_{s,t}(i,j) =
\begin{cases}
\max(i,j), & \text{si } \min(i,j) = 0, \\[10pt]
\min
	\begin{cases}
		L_{s,t}(i-1,j) + 1, \\
		L_{s,t}(i,j-1) + 1, \\
		L_{s,t}(i-1,j-1) + 
	  \begin{cases}
		  1 & \text{si } a_i \neq b_j, \\
		  0 & \text{si } a_i = b_j, \\
	  \end{cases} \\
	\end{cases}
& \text{altrament}.
\end{cases}
\]

\subsection{Centroide d’un clúster}

Donat un clúster
\[
  C_p = \{X_{1}^{(p)}, X_{2}^{(p)}, \dots, X_{N_p}^{(p)}\},
\]
definim el seu centroide com:
\[
  \text{centroide}(C_p)=
  [\bar{X}_{.1}^{(p)}, \bar{X}_{.2}^{(p)}, \dots, \bar{X}_{.N}^{(p)}],
\]
on cada component es calcula segons la tipologia:

\begin{itemize}
  \item \textbf{Numèrica:} mitjana aritmètica.
  \item \textbf{Qualitativa ordenada:} moda ordinal.
  \item \textbf{Qualitativa no ordenada:} moda simple.
  \item \textbf{Multivaluada:} unió/mode d’elements més freqüents.
  \item \textbf{Text lliure (mediana):}
        seleccionem la paraula clau (o seqüència) amb
        \textbf{mínima distància mitjana de Levenshtein} respecte les respostes del clúster.
\end{itemize}

Aquesta “mediana semàntica” és l’anàleg de la mitjana numèrica per strings.

\subsection{Particularitats de la nostra implementació}

La nostra implementació de l’algoritme \textit{k-means} adapta els passos generals al cas de dades d’enquestes de tipologia heterogènia:

\begin{itemize}
  \item \textbf{Inicialització:} seleccionem aleatòriament $k$ individus de les dades com a centroides inicials.
  \item \textbf{Criteri de parada:} l’algoritme s’atura quan:
  \begin{itemize}
    \item cap persona canvia de clúster entre dues iteracions consecutives (convergència), o
    \item s’arriba al nombre màxim d’iteracions establert.
  \end{itemize}
  \item \textbf{Clusters buits en el M-step:} si un clúster queda buit,
        \textbf{assignem com a nou centroide un individu escollit aleatòriament}
        entre tots els disponibles.
  \item Per tal de calcular la distància de Levenshtein entre dos strings $s$ i $t$,
es construeix una taula de programació dinàmica de dimensions 
$(|s|+1) \times (|t|+1)$, on la posició $(i,j)$ conté el valor
$L_{s,t}(i,j)$, és a dir, la distància mínima entre els prefixos
$s[1..i]$ i $t[1..j]$.

La idea fonamental és que qualsevol transformació òptima entre els dos
prefixos ha de derivar d’una de tres operacions bàsiques:

\begin{itemize}
    \item \textbf{Inserció}: transformar $s[1..i]$ en $t[1..j-1]$ i inserir $t_j$,
    amb un cost de $L_{s,t}(i,j-1)+1$.
    \item \textbf{Eliminació}: transformar $s[1..i-1]$ en $t[1..j]$ i eliminar $s_i$,
    amb un cost de $L_{s,t}(i-1,j)+1$.
    \item \textbf{Substitució}: transformar $s[1..i-1]$ en $t[1..j-1]$ i substituir
    $s_i$ per $t_j$, amb un cost de
    \[
    L_{s,t}(i-1,j-1) + 
    \begin{cases}
        0, & \text{si } s_i = t_j,\\
        1, & \text{si } s_i \neq t_j.
    \end{cases}
    \]
\end{itemize}

Així doncs, per a $i,j > 0$, el valor de la taula s’obté aplicant:

\[
L_{s,t}(i,j)=
\min\left\{
\begin{aligned}
    &L_{s,t}(i-1,j)+1,\\
    &L_{s,t}(i,j-1)+1,\\
    &L_{s,t}(i-1,j-1) + 
      \begin{cases}
          0 & \text{si } s_i = t_j,\\
          1 & \text{si } s_i \neq t_j.
      \end{cases}
\end{aligned}
\right\}.
\]

La condició base correspon a les files i columnes inicials. Si un dels
prefixos és buit, la distància coincideix amb la longitud de l’altre prefix:

\[
L_{s,t}(i,0) = i, \qquad
L_{s,t}(0,j) = j.
\]

Un cop omplerta tota la taula, el valor final de la distància de
Levenshtein és $L_{s,t}(|s|,|t|)$.
\end{itemize}

Aquestes modificacions garanteixen l’estabilitat de l’algoritme i eviten fallades en cas de dades molt disperses o de dimensiò alta.

\subsection{Aplicació general del procés}

El flux complet és:

\begin{enumerate}
  \item Inicialitzar aleatòriament els centroides.
  \item Assignar cada persona al clúster amb distància heterogènia mínima.
  \item Calcular els centroides segons les regles de cada tipologia.
  \item Si un clúster està buit, redefinir-ne el centroide aleatòriament.
  \item Repetir fins a convergència o fins al límit d'iteracions.
\end{enumerate}

Aquesta adaptació permet extreure perfils o prototipus de comportament
a partir de respostes d’enquestes complexes i heterogènies.


\end{document}

